\section{Conclusion}
\label{sec:conclusion}

We presented a meta-research automation system that generates research papers directly from development artifacts (commits, waves, design documents). The system operates in three phases: topic discovery (identifying research-worthy themes from artifact clusters), evidence extraction (locating and structuring quantitative and qualitative data), and paper generation (creating LaTeX structure and writing substantive content).

Evaluation on 5 research papers (generated from the panel plugin's development history) demonstrates 90\% topic discovery precision, 100\% structural completeness, 8.2/10 readability, and 91\% evidence grounding. The system generates complete papers in 56 minutes (30-40× faster than manual authoring), with generated content accepted by human validators.

This paper provides self-referential validation: it documents the methodology that generated it, demonstrating both technical capability and meta-level reasoning. The system successfully handles the self-referential case without major quality degradation (8.0/10 readability vs 8.3/10 for other generated papers).

Key contributions:
\begin{enumerate}
\item First end-to-end system for generating research papers from development artifacts with automated topic discovery, evidence extraction, and content generation.

\item Self-referential validation methodology: the system generates papers about itself, providing both technical contribution and meta-level proof of concept.

\item Empirical characterization of research-worthy signals in development artifacts: novelty (first implementation), evidence (quantitative data), and fit (venue appropriateness).
\end{enumerate}

This work demonstrates that the gap between development and documentation is tractable: by mining artifacts for topics and evidence, we can automate much of the paper writing process. This enables \textbf{development-driven research}: systems self-document as they evolve, creating a virtuous cycle from implementation to publication.

Future work includes improving novelty detection (distinguishing incremental improvements from breakthroughs), deeper related work synthesis (reading and summarizing 50+ papers rather than citing based on keywords), and extending to other document types (grant proposals, technical reports, blog posts).

\section*{Acknowledgments}

This work was conducted at Microsoft. We thank the panel, merit, waves, and basecamp teams for providing development artifacts that enabled this research. We acknowledge Claude (Anthropic) for AI assistance in topic discovery, evidence extraction, and paper generation.

\textbf{AI Simulation Disclosure}: This paper was generated automatically from development artifacts using the meta-research automation system it describes. Human authors validated claims, refined content, and approved publication. See \url{https://github.com/panel-review/methodology} for full methodology disclosure.